% ------------------------------------------------------------------------------
\section{Abrir ConcretoLab por primera vez}
% ------------------------------------------------------------------------------

\paso{1}{Abra el programa}

Puede abrirlo de cualquiera de estas formas:

\begin{itemize}
	\item \textbf{Menú Inicio:} pulse la tecla \tecla{Windows} (o el botón Inicio), escriba \textbf{ConcretoLab} y pulse
	      \tecla{Enter}. También está en la lista de aplicaciones, dentro de la carpeta \textbf{ConcretoLab}.
	\item \textbf{Escritorio:} haga doble clic en el acceso \textbf{ConcretoLab}, si lo creó en el paso~5.
\end{itemize}

\begin{consejo}
Para tenerlo siempre a mano, en el menú Inicio haga clic derecho sobre ConcretoLab y elija \menu{Anclar a la barra de
tareas}.
\end{consejo}

\paso{2}{Espere a que cargue}

La primera vez tarda unos segundos más, porque prepara su base de datos. Luego se abre la pantalla de inicio:

\captura[0.86]{guia/11-inicio}{ConcretoLab recién instalado, en la pantalla de inicio.}{fig:primera}

Encima aparece un breve \textbf{recorrido de bienvenida} de cinco páginas que presenta el programa:

\captura[0.8]{guia/10-bienvenida}{Recorrido de bienvenida del primer uso.}{fig:g-bienvenida}

Avance con \boton{Siguiente} o ciérrelo con la \boton{X}. Si marca \textbf{No volver a mostrar al abrir}, no aparecerá
en los próximos inicios (puede volver a verlo desde \menu{Ayuda}).

Arriba están la barra de menús y las pestañas con todas las secciones del programa. Para aprender a usarlas, pulse \menu{Ayuda}
$\rightarrow$ \menu{Manual de usuario}: el manual se abre dentro del programa. En cada pantalla, el botón \boton{?} de
arriba a la derecha (o la tecla \tecla{F1}) lo abre justo en la sección que necesita. Para practicar sin tocar sus
datos, use el botón \boton{Agregar proyecto de ejemplo} de la pantalla de inicio.

\paso{3}{Solo si usaba una versión anterior con contraseña}

Las versiones anteriores de ConcretoLab pedían una contraseña para abrir los datos. \textbf{La versión actual abre
directamente, sin contraseña.} Si en esta computadora ya había datos protegidos con contraseña, la primera vez aparecerá
la ventana \textbf{«Abrir datos de una versión anterior»}:

\begin{enumerate}
	\item Escriba la contraseña que usaba antes y pulse \boton{Abrir}.
	\item Aparecerá el mensaje \textbf{«Datos listos»}: sus proyectos se conservaron y desde ese momento ConcretoLab ya
	      no pedirá contraseña.
\end{enumerate}

\begin{nota}
Antes de quitar la contraseña, ConcretoLab guarda una copia exacta de la base anterior. La verá en la sección
\menu{Datos y copias} con el tipo «Base anterior con contraseña». Si se equivoca de contraseña, tiene tres intentos y
sus datos no se modifican; si pulsa \boton{Salir}, el programa se cierra sin tocar nada.
\end{nota}

\subsection{Dónde se guardan sus datos}

\begin{table}[H]
	\centering
	\caption{Carpetas que usa ConcretoLab.}
	\label{tab:carpetas}
	\begin{tabular}{L{4.6cm}L{9.2cm}}
		\hline
		\textbf{Qué} & \textbf{Dónde} \\ \hline
		El programa & \ruta{C:\textbackslash Users\textbackslash <su usuario>\textbackslash AppData\textbackslash Local\textbackslash ConcretoLab} \\
		Sus proyectos (base de datos) & \ruta{C:\textbackslash Users\textbackslash <su usuario>\textbackslash .concretolab\textbackslash concretolab.db} \\
		Copias de seguridad & \ruta{C:\textbackslash Users\textbackslash <su usuario>\textbackslash .concretolab\textbackslash backups} \\
		Registros técnicos & \ruta{C:\textbackslash Users\textbackslash <su usuario>\textbackslash .concretolab\textbackslash logs} \\ \hline
	\end{tabular}
\end{table}

No necesita abrir estas carpetas: todo se maneja desde \menu{Datos y copias}, dentro del programa. Allí también puede
crear una copia de seguridad cuando quiera y llevar un proyecto a otra computadora.
